

% This is a simple sample document.  For more complicated documents take a look in the exercise tab. Note that everything that comes after a % symbol is treated as comment and ignored when the code is compiled.

\documentclass[fleqn]{report} % \documentclass{} is the first command in any LaTeX code.  It is used to define what kind of document you are creating such as an article or a book, and begins the document preamble

\usepackage{amsmath} % \usepackage is a command that allows you to add functionality to your LaTeX code
\usepackage{amsfonts}
\usepackage{algorithm}
\usepackage{algpseudocode}
\usepackage{graphicx}
\usepackage{xcolor}
\newcommand\ddfrac[2]{\frac{\displaystyle #1}{\displaystyle #2}}
%\usepackage[linesnumbered,ruled,vlined]{algorithm2e}



\title{Technical report\\Sparsifying $l_2$ regularized least squares regression by biconjugate convex relaxation} % Sets article title
\author{Dr. Mamadou THIAO} % Sets authors name
\date{\today} % Sets date for date compiled
\newtheorem{thm}{Proposition}[subsection]
\newenvironment{proof}{\paragraph{Proof:}}{\hfill$\square$}
\newtheorem{crly}{Corollary}[subsection]

% The preamble ends with the command \begin{document}
\begin{document} % All begin commands must be paired with an end command somewhere
    \maketitle % creates title using information in preamble (title, author, date)
    \tableofcontents
    \section{Linear least squares} % creates a section
\begin{equation}
		\label{Ridge_Regularization}
		(\text{RIDGE}^{\beta})\;
		\left\{
		\begin{array}{ll}\min & \left\|Ax-b\right\|^2_2
			+ \beta\left\|x\right\|_2^2\\
			s.t.&\;x\in\mathbb{R}^n,\\
		\end{array} 
		\right.
    \end{equation}
		where $\beta> 0$ is a parameter.
    \section{Sparsifying linear least squares} % creates a section
We use a $l_0$ regularization approach to promote sparse solutions as expressed in the following problem
		\begin{equation}
		\left\{
		\begin{array}{ll}\min & \left\|Ax-b\right\|^2_2 + \rho\left\|x\right\|_0 + \beta\left\|x\right\|^2_2  \\
			s.t.&\;x\in\mathbb{R}^n.\\
		\end{array} 
		\right.
    \end{equation}
		where $\rho \geq 0$ is a parameter which controls the sparsity of the solutions. \\
With $\delta := \sqrt{\ddfrac{\rho}{\beta}}$, we have $\rho = \beta\delta^2$. The previous problem  becomes
\begin{equation}
		\label{L0_L2_Regularized_Sparse_Approximation}
		(\mathbb{P}^{\beta, \delta})\;
		\left\{
		\begin{array}{ll}\min & \left\|Ax-b\right\|^2_2 + \beta\delta^2\left\|x\right\|_0 + \beta\left\|x\right\|^2_2 \\
			s.t.&\;x\in\mathbb{R}^n.\\
		\end{array} 
		\right.
    \end{equation}

		\section{Relaxation $(\mathbb{Q}^{\beta,\delta})$}
Replacing $\left\|.\right\|_0 +\frac{1}{\delta^2} \left\|.\right\|^2_2$ by its tightest convex under approximation, its biconjugate, 
		\[
			\left(\left\|.\right\|_0 + \frac{1}{\delta^2} \left\|.\right\|_2^2\right)^{**}(x),
		\]
		see Le Thi et al., we obtain the following biconjugate relaxation (\textbf{B}i\textbf{C}onjugate \textbf{R}elaxation,  \textbf{BCR} )
		\[
		(\mathbb{Q}^{\beta,\delta})\left\{\begin{array}{ll}\min & F(x):=\left\|Ax-b\right\|^2_2 + \beta\delta^2\left(\left\|.\right\|_0 + \frac{1}{\delta^2}\left\|.\right\|_2^2\right)^{**}(x)\\
		s.t.&\;x\in\mathbb{R}^n.
		\end{array} \right.
		\]
$(\mathbb{Q}^{\beta,\delta})$ is a convex under approximation of $(\mathbb{P}^{\beta,\delta})$ and $(\mathbb{Q}^{\beta,\delta})$ solves $(\mathbb{P}^{\beta,\delta})$ when its solution satisfies particular global optimality conditions (see Le Thi et al.).
		\[
			\left(\left\|.\right\|_0 + \ddfrac{1}{\delta^2}\left\|.\right\|_2^2\right)^{**}(x) = \sum_{i}^{n}\left(\left|.\right|_0 + \ddfrac{1}{\delta^2} \left|.\right|_2^2\right)^{**}(x_i),
		\]
\[
\left(\left|.\right|_0 + \ddfrac{1}{\delta^2} \left|.\right|_2^2\right)^{**}(x_i)=\ddfrac{1}{\delta^2} x_i^2  - \ddfrac{1}{\delta^2} \left[\left(\delta-\left|x_i\right|\right)^+\right]^2 +1=\left\{ \begin{array}{ll}
 \ddfrac{1}{\delta^2}x_i^2 + 1 &\;\text{if}\;\left|x_i\right|\geq \delta.\\
&\\
 2\ddfrac{1}{\delta}\left|x_i\right|& \text{otherwise}.
\end{array}
\right.
\]
\subsection{Optimality conditions for $(\mathbb{Q}^{\beta,\delta})$}
The following  form for the optimality conditions are very useful for the theoretical study of the links between the biconjugate relaxation BCR, Ridge, Lasso and ElasticNet.
\subsubsection{Optimality condition for $(\mathbb{Q}^{\beta,\delta})$}
The optimality condition for the convex relaxation $(\mathbb{Q}^{\beta,\delta})$ is given by, Le Thi et al.
\[
	-2A^T(Ax-b)\in \beta\delta^2\partial \left(\left\|.\right\|_0 +  \ddfrac{1}{\delta^2} \left\|.\right\|_2^2\right)^{**}(x).
\]
As
\[
\partial\left(\left|.\right|_0 + \ddfrac{1}{\delta^2} \left|.\right|_2^2\right)^{**}(x_i)=\left\{ \begin{array}{ll}
\Bigg\{2\ddfrac{1}{\delta^2}  x_i\Bigg\}&\;\text{if}\; \left|x_i\right|\geq \delta,\\
&\\
\Bigg\{2\ddfrac{1}{\delta} \;\text{sign}(x_i)\Bigg\} &\;\text{if}\; 0<\left|x_i\right|\leq \delta,\\
&\\
\Bigg[-2\ddfrac{1}{\delta} , 2\ddfrac{1}{\delta} \Bigg]&\;\text{if}\; x_i=0,
\end{array}
\right.
\]
the condition becomes
\[
\left(\underset{opt}{\mathbb{Q}^{\beta,\delta}}\right)\;
\left[A^T(Ax-b)\right]_i\in 
\left\{ \begin{array}{ll}
\left\{-\beta x_i\right\}&\;\text{if}\; \left|x_i\right|\geq\delta\\
\left\{-\delta\beta\;\text{sign}(x_i)\right\} &\;\text{if}\; 0<\left|x_i\right|\leq \delta\\
\left[-\delta\beta, \delta\beta\right]&\;\text{if}\; x_i=0.
\end{array}
\right.
\]
or
\[
\left[A^T(Ax-b) +\beta x\right]_i\in 
\left\{ \begin{array}{ll}
\left\{0\right\}&\;\text{if}\; \left|x_i\right|\geq\delta\\
\left\{\beta x_i-\delta\beta\;\text{sign}(x_i)\right\} &\;\text{if}\; 0<\left|x_i\right|\leq \delta\\
\left[-\delta\beta, \delta\beta\right]&\;\text{if}\; x_i=0.
\end{array}
\right.
\]
\begin{thm}
$x=0$ will be an optimal solution for $(\mathbb{Q}^{\beta,\delta})$ if and only if $\left\|A^Tb\right\|_{\infty}\leq\delta\beta $.
\end{thm}
\begin{proof}
    Direct use of $\left(\underset{opt}{\mathbb{Q}^{\beta,\delta}}\right)$.
\end{proof}
\begin{thm}If  $x$ is a solution for $(\mathbb{Q}^{\beta,\delta})$ and $\delta<
\underset{i\in \text{supp}(x)}{\min}\left|x_i\right|$,\\ then $x$ is a solution for $(\mathbb{Q}^{\beta,\delta^{\prime}})$, for all $\delta^{\prime}$, $\delta\leq \delta^{\prime}\leq \underset{i\in \text{supp}(x)}{\min}\left|x_i\right|$.
\end{thm}
\begin{proof}
    Direct use of $\left(\underset{opt}{\mathbb{Q}^{\beta,\delta}}\right)$.
\end{proof}
\subsection{Dual of $(\mathbb{Q}^{\beta,\delta})$}
Lagrangian dual of $(\mathbb{D}^{\beta,\delta})$ is 
\[
		(\mathbb{D}^{\beta,\delta})\;\left\{\begin{array}{ll}
\max & G(\nu):=-\ddfrac{1}{4}\left\|\nu\right\|^2_2 - b^T\nu - \beta\delta^2\left(\left\|.\right\|_0 + \ddfrac{1}{\delta^2}  \left\|.\right\|_2^2\right)^{*}(-\ddfrac{A^T\nu}{\beta\delta^2})\\
		s.t.&\;\nu\in\mathbb{R}^m.
		\end{array} \right.
		\]
With 
		\[
			\left(\left\|.\right\|_0 + \ddfrac{1}{\delta^2}  \left\|.\right\|_2^2\right)^{*}(z) = \sum_{i}^{m}\left(\left|.\right|_0 + \ddfrac{1}{\delta^2} \left|.\right|_2^2\right)^{*}(z_i),
		\]
\[
			\left(\left|.\right|_0 + \ddfrac{1}{\delta^2}  \left|.\right|_2^2\right)^{*}(z_i) = \left(\ddfrac{\delta^2}{4}z_i^2 - 1\right)^+=\left\{ \begin{array}{ll}
\ddfrac{\delta^2}{4}z_i^2 - 1 &\;\text{if}\; \left|z_i\right|\geq \ddfrac{2}{\delta}\\
0& \text{otherwise}.
\end{array}
\right.
		\]

Weak duality holds. As Slater's conditition is satisfied, strong duality holds for $(\mathbb{D}^{\beta,\delta})$.
\subsubsection{Optimality condition for $(\mathbb{D}^{\beta,\delta})$}
The optimality condition for the convex program $(\mathbb{D}^{\beta,\delta})$ is given by,
\[
	\ddfrac{1}{2}\nu+b\in A\partial (\left\|.\right\|_0 + \ddfrac{1}{\delta^2}  \left\|.\right\|_2^2)^{*}(-\ddfrac{A^T\nu}{\beta\delta^2}).
\]
As
\[
\partial (\left\|.\right\|_0 + \ddfrac{1}{\delta^2}  \left\|.\right\|_2^2)^{*}(z_i)=\left\{ \begin{array}{ll}
\left\{\ddfrac{\delta^2}{2}z_i\right\}&\;\text{if}\; \left|z_i\right|> \ddfrac{2}{\delta}  \\
\left\{0\right\} &\;\text{if}\; \left|z_i\right|<  \ddfrac{2}{\delta} \\
\left[0, \delta\right]&\;\text{if}\; z_i=  \ddfrac{2}{\delta} ,\\
\left[-\delta, 0\right]&\;\text{if}\; z_i= - \ddfrac{2}{\delta} ,
\end{array}
\right.
\]
the condition becomes
\[
\left(\underset{opt}{\mathbb{D}^{\beta,\delta}}\right)\;
\ddfrac{1}{2}\nu+b = Ay,\; 
y_i\in 
\left\{ \begin{array}{ll}
\left\{-\ddfrac{1}{2\beta}\left[A^T\nu\right]_i\right\}&\;\text{if}\; \left|\left[A^T\nu\right]_i\right|> 2\delta\beta\\
\left\{0\right\} &\;\text{if}\; \left|\left[A^T\nu\right]_i\right|< 2\delta\beta\\
\left[0, \delta\right]&\;\text{if}\; \left[A^T\nu\right]_i= -2\delta\beta,\\
\left[-\delta, 0\right]&\;\text{if}\; \left[A^T\nu\right]_i= 2\delta\beta.
\end{array}
\right.
\]
\subsubsection{Finding a dual feasible from a primal feasible}
The previous primal and dual optimality conditions show that if $x$ is a primal solution then $\nu=Ax-b$ is a dual solution.\\ 
Therefore for an arbitray feasible primal point $x$ we define a dual feasible point $\nu = s\nu_0$,  $\nu_0 = (Ax-b)$. The scalar $s$ will be the argmax of the following 1D convex program:
\[
		\left\{\begin{array}{ll}
\max & -\ddfrac{1}{4}\left\|\nu_0\right\|^2_2 s^2 - b^T\nu_0 s - \beta\delta^2\left(\left\|.\right\|_0 + \ddfrac{1}{\delta^2} \left\|.\right\|_2^2\right)^{*}(-\ddfrac{A^T\nu_0}{\beta\delta^2}s)\\
		s.t.&\;s\in\mathbb{R},
		\end{array} \right.
		\]
say 
\[
		\left\{\begin{array}{ll}
\max & -\ddfrac{1}{4}\left\|\nu_0\right\|^2_2 s^2 - b^T\nu_0 s - \beta\delta^2\sum_{i=1}^{n} \left(\ddfrac{\left[A^T\nu_0\right]_i^2}{4\delta^2\beta^2}s^2 - 1\right)^+\\
		s.t.&\;s\in\mathbb{R}.
		\end{array} \right.
		\]
This is a piecewise 1D quadratic convex program. A solution will have necessarily the opposite sign of $b^T\nu_0$.\\
 We will first solve 
\[
		\left\{\begin{array}{ll}
\max & -\ddfrac{1}{4}\left\|\nu_0\right\|^2_2 s^2 + \left|b^T\nu_0\right| s -  \beta\delta^2\sum_{i=1}^{n} \left(\ddfrac{\left[A^T\nu_0\right]_i^2}{4\delta^2\beta^2}s^2 - 1\right)^+\\
		s.t.&\;s\geq0,
		\end{array} \right.
		\]
and then update the sign of $s$.
\subsubsection{Duality gap}
$G(\nu)$ lower bounds $F(x)$ and the corresponding duality gap
\[
\eta := F(x) - G(\nu)
\]
is always nonnegative from weak duality. $x$ is no more than $\eta$-suboptimal. At the optimum the gap is zero from strong duality.
\subsection{Theoretical links with Ridge and Lasso}
\subsubsection{Theoretical links with Ridge}
Recall
\begin{equation}
		(\text{RIDGE}^{\beta})\;
		\left\{
		\begin{array}{ll}\min & \left\|Ax-b\right\|^2_2 + \beta\left\|x\right\|_2^2\\
			s.t.&\;x\in\mathbb{R}^n.\\
		\end{array} 
		\right.
    \end{equation}
Ridge optimality condition ($2A^T(Ax-b) + 2\beta x = 0$) can be written as
\[
\left(\underset{opt}{\text{RIDGE}^{\beta}}\right)
\left[A^T(Ax-b)\right]_i\in 
\left\{ 
-\beta x_i\right\}.
\]
\begin{thm}A $\text{RIDGE}$ is always a $(\mathbb{Q}^{\beta,\delta})$.\\\\
More precisely, assume $A$ has no null column,
\begin{enumerate}
\item If $x$ is a solution for $(\text{RIDGE}^{\beta})$,\\ then $x$ is a solution for $(\mathbb{Q}^{\beta,\delta})$ for all $\delta\leq \underset{i}{\min}\left|x_i\right|$.
\item If $x$ is a solution for $(\mathbb{Q}^{\beta,\delta})$ and $\delta\leq \underset{i}{\min}\left|x_i\right|$,\\ then $x$ is a solution for $(\text{RIDGE}^{\beta})$.
\end{enumerate}
\end{thm}
\begin{proof}
    Direct use of $\left(\underset{opt}{\text{RIDGE}^{\beta}}\right)$ and $\left(\underset{opt}{\mathbb{Q}^{\beta,\delta}}\right)$.
\end{proof}
\subsubsection{Theoretical links with Lasso}
\begin{equation}
		\label{Lasso_Regularization}
		(\text{LASSO}^{\gamma})\;
		\left\{
		\begin{array}{ll}\min & \left\|Ax-b\right\|^2_2 + \gamma\left\|x\right\|_1\\
			s.t.&\;x\in\mathbb{R}^n.\\
		\end{array} 
		\right.
    \end{equation}
Lasso optimality condition can be written as
\[
\left(\underset{opt}{\text{LASSO}^{\gamma}}\right)\;\left[A^T(Ax-b)\right]_i\in 
\left\{ \begin{array}{ll}
\left\{-\ddfrac{\gamma}{2}\;\text{sign}(x_i)\right\} &\;\text{if}\; x_i\neq 0\\
\left[-\ddfrac{\gamma}{2}, \ddfrac{\gamma}{2}\right]&\;\text{if}\; x_i =0.
\end{array}
\right.
\]
\begin{thm} A $\text{LASSO}$ is always a $(\mathbb{Q}^{\beta,\delta})$.\\\\
More precisely,
\begin{enumerate}
\item If $x$ is a solution for $(\text{LASSO}^{\gamma})$,\\ then  $x$ is a solution for $(\mathbb{Q}^{\ddfrac{\gamma}{2\delta},\delta})$ for all $\delta\geq
\underset{i\in\text{supp}(x)}{\max}\left|x_i\right|$.
\item If $x$ is a solution for $(\mathbb{Q}^{\beta,\delta})$ and $\delta\geq\underset{i\in\text{supp}(x)}{\max}\left|x_i\right|$, \\ then $x$ is a solution for $(\text{LASSO}^{2\delta\beta})$.
\end{enumerate}
\end{thm}
\begin{proof}
    Direct use of $\left(\underset{opt}{\text{LASSO}^{\gamma}}\right)$ and $\left(\underset{opt}{\mathbb{Q}^{\beta,\delta}}\right)$.
\end{proof}
\subsubsection{From Ridge to Lasso. How navigate $(\mathbb{Q}^{\beta,\delta})$}
Consider $(\text{RIDGE}^{\beta})$, for a $\beta>0$. (Assume $A$ has no null column)\\ Define 
\[
\delta_{\text{RIDGE}^{\beta}}:=\underset{i\in\text{supp}(x^{\text{RIDGE}^{\beta}})}{\min} \left|x^{\text{RIDGE}^{\beta}}_i\right|.
\]
Let us analyze $\delta$s values.\\
\\
We already have $x^{\text{RIDGE}^{\beta}}$ solution of  $(\mathbb{Q}^{\beta,\delta})$ for all $\delta\leq\delta_{\text{RIDGE}^{\beta}}$.\\\\
We will consider $\delta_{\text{RIDGE}^{\beta}} > \delta$.\\
\\
From inequality $F(x^{\mathbb{Q}^{\beta,\delta}}) \leq F(0)$, we deduce
\[
  \left|x^{\mathbb{Q}^{\beta,\delta}}_i\right|\leq\left\|x^{\mathbb{Q}^{\beta,\delta}}\right\|_{1}  \leq\ddfrac{1}{2\delta\beta}F(0),\forall\;i.
\]
So it suffices to constrain $\ddfrac{1}{2\delta\beta}F(0)\leq \delta$ to ensure $\left|x^{\mathbb{Q}^{\beta,\delta}}_i\right|\leq \delta,$ for all $i$.\\\\
So $x^{\mathbb{Q}^{\beta,\delta}}$ is a solution of $( \text{LASSO}^{2\delta\beta})$,  for all $\delta > \max\left\{\sqrt{\ddfrac{1}{2\beta}F(0)},  \delta_{\text{RIDGE}^{\beta}}\right\}$.\\
\\
Define 
\[
\delta_{\text{RIDGE}^{\beta},\text{LASSO}} :=\inf\left\{\delta  >  \delta_{\text{RIDGE}^{\beta}}:\;x^{\mathbb{Q}^{\beta,\delta}}\; \text{is a solution of }( \text{LASSO}^{2\delta\beta})\right\}.
\]
We also have necessarily  
\[\rho_{\text{RIDGE}^{\beta},\text{LASSO}}\leq\max\left\{\sqrt{\ddfrac{1}{2\beta}F(0)},  \delta_{\text{RIDGE}^{\beta}}\right\}.
\]
\begin{thm} 
\begin{enumerate}
\item When $\delta$ moves from $ \delta_{\text{RIDGE}^{\beta}}$ to $\delta_{\text{RIDGE}^{\beta},\text{LASSO}}$ , $(\mathbb{Q}^{\beta,\delta})$ moves from $\text{RIDGE}$ to $\text{LASSO}$.
\begin{itemize}
\item before $ \delta_{\text{RIDGE}^{\beta}}$, $(\mathbb{Q}^{\beta,\delta})$ is a $\text{RIDGE}$,
\item after $ \delta_{\text{RIDGE}^{\beta},\text{LASSO}}$, $(\mathbb{Q}^{\beta,\delta})$ is a $\text{LASSO}$.
\end{itemize}
\item When $\delta$ moves strictly between $ \delta_{\text{RIDGE}^{\beta}}$ and $\delta_{\text{RIDGE}^{\beta},\text{LASSO}}$ , $(\mathbb{Q}^{\beta,\delta})$ cannot be a $\text{RIDGE}$ nor a $\text{LASSO}$.\\
The components of $x^{\mathbb{Q}^{\beta,\delta} }$ satisfy $\left(\underset{opt}{\mathbb{Q}^{\beta,\delta}}\right)$
\[
\left[A^T(Ax^{\mathbb{Q}^{\beta,\delta}}-b)\right]_i\in 
\left\{ \begin{array}{ll}
\left\{-\beta x^{\mathbb{Q}^{\beta,\delta}}_i\right\}&\;\text{if}\; \left|x^{\mathbb{Q}^{\beta,\delta}}_i\right|\geq \delta\\
\left\{-\delta\beta\;\text{sign}(x^{\mathbb{Q}^{\beta,\delta}}_i)\right\} &\;\text{if}\; 0<\left|x^{\mathbb{Q}^{\beta,\delta}}_i\right|\leq \delta\\
\left[-\delta\beta, \delta\beta\right]&\;\text{if}\; x^{\mathbb{Q}^{\beta,\delta}}_i=0.
\end{array}
\right.
\]
\end{enumerate}
\end{thm}
The second point of the previous proposition highlights the mixing role (on the components of $x$) played by $(\mathbb{Q}^{\beta,\delta} )$. The first line of the optimality condition is for RIDGE, with $\beta$, and the followings are for LASSO,  with $\gamma:=2\delta\beta$. The value $\delta=\ddfrac{\gamma}{2\beta}$ defines a trade-off line on the coordinates value.
\subsection{Mixing roles: $(\mathbb{Q}^{\beta,\delta} )$  vs $(\text{ELASTICNET}^{\beta^{\prime},\gamma})$}
\subsubsection{Theoretical links with ElasticNet}
\begin{equation}
		\label{ElasticNet_Regularization}
		(\text{ELASTICNET}^{\beta^{\prime},\gamma})\;
		\left\{
		\begin{array}{ll}\min & \left\|Ax-b\right\|^2_2 + \beta^{\prime}\left\|x\right\|_2^2 + \gamma\left\|x\right\|_1\\
			s.t.&\;x\in\mathbb{R}^n.\\
		\end{array} 
		\right.
    \end{equation}
ElasticNet optimality condition can be written as
\[
\left(\underset{opt}{\text{ELASTICNET}^{\beta^{\prime},\gamma}}\right)\;\left[2A^T(Ax-b)\right]_i\in 
\left\{ \begin{array}{ll}
\left\{-2\beta^{\prime} x_i -\gamma\;\text{sign}(x_i)\right\} &\;\text{if}\; x_i\neq 0\\
\left[-\gamma,  \gamma\right]&\;\text{if}\; x_i =0.
\end{array}
\right.
\]
\begin{thm} $x$ is a common solution for $(\text{ELASTICNET}^{\beta^{\prime},\gamma})$ and $(\mathbb{Q}^{\beta,\delta} )$ if and only if 
$x$ is a solution of $(\text{ELASTICNET}^{\beta^{\prime},\gamma})$ or $(\mathbb{Q}^{\beta,\delta} )$ and 
one of the following conditions holds:
\begin{enumerate}
    \item there exists $C_1$, such that
\[\left|x_i\right|=C_1,\forall\;i,\;x_i\neq0,
\]
\[C_1 \geq \delta,\;\beta=\beta^{\prime} + \ddfrac{\gamma}{2C_1}.
\]
    \item there exists $C_2$, such that
\[\left|x_i\right|=C_2,\forall\;i,\;x_i\neq0,
\]
\[C_2 \leq \delta,\;\delta\beta=\ddfrac{\gamma}{2}+\beta^{\prime} C_2.
\]
    \item there exist $C_1$ and $C_2$, $C_2 \leq C_1$, such that
\[\left|x_i\right|=C_1\;\text{or}\;C_2,\forall\;i,\;x_i\neq0,
\]
\[ \delta = \ddfrac{C_1(\gamma+2\beta^{\prime} C_2)}{\gamma+2\beta^{\prime} C_1},\;\delta\beta=\ddfrac{\gamma}{2}+\beta^{\prime} C_2.
\]
\end{enumerate}\end{thm}
\begin{proof}
    Use $\left(\underset{opt}{\text{ELASTICNET}^{\beta^{\prime},\gamma}}\right)$,  $\left(\underset{opt}{\mathbb{Q}^{\beta,\delta}} \right)$ and compare solutions.
\end{proof}
\section{Solving $(\mathbb{Q}^{\beta,\delta})$}
Let $\beta>0$ and let us denote by $x^{\beta}(\delta)$ a solution for $(\mathbb{Q}^{\beta,\delta})$. The solution set with respect to $\delta$ is called the path and can be expressed as 
\[
	\textbf{PATH}^{\beta}:= \left\{x^{\beta}(\delta):\;\delta\geq 0\right\}.
\] 
\subsubsection{Computing a new solution on the path from a known one}
Here we will consider that we have $(\mathbb{Q}^{\beta,\underline{\delta}})$, $\underline{\delta}>0$ and a corresponding solution $\underline{x}:=x^{\beta}(\underline{\delta})$.\\
\\
Let us define the following partition of the indices $i$, $i=1,..., n$.
\[
\begin{array}{l}
\underline{\mathbb{K}}_0:=\left\{i:\;\left|\underline{x}_i\right|>\underline{\delta} \right\},\\
\underline{\mathbb{K}}_1\cup\underline{\mathbb{K}}_2:=\left\{i:\;\left|\underline{x}_i\right|=\underline{\delta} \right\},\; \underline{\mathbb{K}}_1\cap\underline{\mathbb{K}}_2=\emptyset,\\
\underline{\mathbb{K}}_3:=\left\{i:\;0<\left|\underline{x}_i\right|<\underline{\delta} \right\},\\
\underline{\mathbb{K}}_4\cup\underline{\mathbb{K}}_5:=\left\{i:\;\underline{x}_i=0,\;\left|\left[A^T(A\underline{x} -b)\right]_i\right|=\underline{\delta} \beta\right\},\;\underline{\mathbb{K}}_4\cap\underline{\mathbb{K}}_5=\emptyset,\\
\underline{\mathbb{K}}_6:=\left\{i:\;\underline{x}_i=0,\;0<\left|\left[A^T(A\underline{x} -b)\right]_i\right|<\underline{\delta} \beta \right\},\\
\underline{\mathbb{K}}_7:=\left\{i:\;\underline{x}_i=0,\;\left[A^T(A\underline{x} -b)\right]_i=0\right\},\\
\underline{M}:=\text{diag}(\underline{\theta}),\\
\underline{\theta}_i:=\text{sign}(\underline{x}_i):=
\left\{
\begin{array}{ll}
+1& \text{if}\;\underline{x}_i >0,\\
- 1& \text{if}\;\underline{x}_i <0,\\
- \text{sign}(\left[A^T(A\underline{x} -b)\right]_i)& \text{if}\;\underline{x}_i=0,\;\left[A^T(A\underline{x} -b)\right]_i\neq0,\\
+1& \text{otherwise}.
\end{array}
\right.
\end{array} 
\]
Let $\tau\in\left\{-1, +1\right\}$ given and consider the following polyhedra
\[
\begin{array}{lrrrrll}
(A^TA\underline{M})_iu+\beta\underline{\theta}_i\textcolor{green}{u_i}&&&&&=(A^Tb)_i&i\in\underline{\mathbb{K}}_0\\
(A^TA\underline{M})_iu+\beta\underline{\theta}_i\textcolor{green}{u_i}&&&&&=(A^Tb)_i&i\in\underline{\mathbb{K}}_1\\
(A^TA\underline{M})_iu+\beta\underline{\theta}_i\textcolor{green}{u_i}&&&+\beta\underline{\theta}_it_i&&=(A^Tb)_i&i\in\underline{\mathbb{K}}_2\\
(A^TA\underline{M})_iu+\beta\underline{\theta}_i\textcolor{green}{u_i}&&&+\beta\underline{\theta}_it_i&&=(A^Tb)_i&i\in\underline{\mathbb{K}}_3\\
(A^TA\underline{M})_iu+\beta\underline{\theta}_i\textcolor{red}{u_i}&&&+\beta\underline{\theta}_it_i&&=(A^Tb)_i&i\in\underline{\mathbb{K}}_4\\
(A^TA\underline{M})_iu&&+\beta\underline{\theta}_i\textcolor{green}{s_i}&&&=(A^Tb)_i&i\in\underline{\mathbb{K}}_5\\
(A^TA\underline{M})_iu&&+\beta\underline{\theta}_i\textcolor{green}{s_i}&&&=(A^Tb)_i&i\in\underline{\mathbb{K}}_6\\
(A^TA\underline{M})_iu&&-\beta\underline{\theta}_i\textcolor{green}{s_i}&&+\beta\underline{\theta}_iw_i&=(A^Tb)_i&i\in\underline{\mathbb{K}}_7\\
u_i&+\tau\gamma&&-\textcolor{green}{t_i}&&=\underline{\delta}&i\in\underline{\mathbb{K}}_0\\
u_i&+\tau\gamma&&-\textcolor{red}{t_i}&&=\underline{\delta}&i\in\underline{\mathbb{K}}_1\\
u_i&+\tau\gamma&&+\textcolor{red}{t_i}&&=\underline{\delta}&i\in\underline{\mathbb{K}}_2\\
u_i&+\tau\gamma&&+\textcolor{green}{t_i}&&=\underline{\delta}&i\in\underline{\mathbb{K}}_3\\
u_i&+\tau\gamma&&+\textcolor{green}{t_i}&&=\underline{\delta}&i\in\underline{\mathbb{K}}_4\\
&\tau\gamma&+s_i&&+\textcolor{red}{w_i}&=\underline{\delta}&i\in\underline{\mathbb{K}}_5\\
&\tau\gamma&+s_i&&+\textcolor{green}{w_i}&=\underline{\delta}&i\in\underline{\mathbb{K}}_6\\
&\tau\gamma&+s_i&&+\textcolor{green}{w_i}&=\underline{\delta}&i\in\underline{\mathbb{K}}_7\\
u_i&&&&&=0&i\in\underline{\mathbb{K}}_5\\
u_i&&&&&=0&i\in\underline{\mathbb{K}}_6\\
u_i&&&&&=0&i\in\underline{\mathbb{K}}_7\\
&&&&&\\
u\geq0&\gamma\geq0&s\geq0&t\geq0&w\geq0&.\\
\end{array} 
\]
The polyhedra contains $u:=\left|\underline{x}\right|$, $\gamma:=0$, ... so is nonempty.
\begin{thm} $x:=\underline{M}u$ satisfies $(\mathbb{Q}^{\beta,\underline{\delta}-\tau\gamma})$ optimality conditions.
\end{thm}
\begin{proof} With $x:=\underline{M}u$, and $\delta:=\underline{\delta}-\tau\gamma$, we have 
\[
\begin{array}{ll}
\left[A^T(Ax -b)+\beta x\right]_i &= \left\{
\begin{array}{ll}
0&i\in\underline{\mathbb{K}}_0\\
0&i\in\underline{\mathbb{K}}_1\\
-\beta\underline{\theta}_it_i&i\in\underline{\mathbb{K}}_2\\
-\beta\underline{\theta}_it_i&i\in\underline{\mathbb{K}}_3\\
-\beta\underline{\theta}_it_i&i\in\underline{\mathbb{K}}_4\\
-\beta\underline{\theta}_is_i&i\in\underline{\mathbb{K}}_5\\
-\beta\underline{\theta}_is_i&i\in\underline{\mathbb{K}}_6\\
\beta\underline{\theta}_is_i-\beta\underline{\theta}_iw_i&i\in\underline{\mathbb{K}}_7\\
\end{array} 
\right.\\
&\\
& = \left\{
\begin{array}{lll}
0&i\in\underline{\mathbb{K}}_0\\
0&i\in\underline{\mathbb{K}}_1\\
\beta x - \delta\beta\text{sign}(x_i)&i\in\underline{\mathbb{K}}_2\\
\beta x - \delta\beta\text{sign}(x_i)&i\in\underline{\mathbb{K}}_3\\
\beta x - \delta\beta\text{sign}(x_i)&i\in\underline{\mathbb{K}}_4\\
-(\delta\beta - w_i\beta)\text{sign}(x_i)&i\in\underline{\mathbb{K}}_5\\
-(\delta\beta - w_i\beta)\text{sign}(x_i)&i\in\underline{\mathbb{K}}_6\\
-(-\delta\beta + 2w_i\beta)\text{sign}(x_i)&i\in\underline{\mathbb{K}}_7\\
\end{array} 
\right.
\end{array} 
\]
Then apply polyhedra constraints on the previous expression.
\end{proof}
\subsubsection{Escaping point with $\gamma=0$ with simplex method.}
Our goal is to maximize $\left|\delta -\underline{\delta}\right|=\left|-\tau\gamma\right|=\gamma$ while keeping the polyhedral constraints valid and  starting with green and red ones as basis.
\[
\begin{array}{lrrrrll}
\text{max}\;\gamma&&&&&&\\
&&&&&&\\
(A^TA\underline{M})_iu+\beta\underline{\theta}_i\textcolor{green}{u_i}&&&&&=(A^Tb)_i&i\in\underline{\mathbb{K}}_0\\
(A^TA\underline{M})_iu+\beta\underline{\theta}_i\textcolor{green}{u_i}&&&&&=(A^Tb)_i&i\in\underline{\mathbb{K}}_1\\
(A^TA\underline{M})_iu+\beta\underline{\theta}_i\textcolor{green}{u_i}&&&+\beta\underline{\theta}_it_i&&=(A^Tb)_i&i\in\underline{\mathbb{K}}_2\\
(A^TA\underline{M})_iu+\beta\underline{\theta}_i\textcolor{green}{u_i}&&&+\beta\underline{\theta}_it_i&&=(A^Tb)_i&i\in\underline{\mathbb{K}}_3\\
(A^TA\underline{M})_iu+\beta\underline{\theta}_i\textcolor{red}{u_i}&&&+\beta\underline{\theta}_it_i&&=(A^Tb)_i&i\in\underline{\mathbb{K}}_4\\
(A^TA\underline{M})_iu&&+\beta\underline{\theta}_i\textcolor{green}{s_i}&&&=(A^Tb)_i&i\in\underline{\mathbb{K}}_5\\
(A^TA\underline{M})_iu&&+\beta\underline{\theta}_i\textcolor{green}{s_i}&&&=(A^Tb)_i&i\in\underline{\mathbb{K}}_6\\
(A^TA\underline{M})_iu&&-\beta\underline{\theta}_i\textcolor{green}{s_i}&&+\beta\underline{\theta}_iw_i&=(A^Tb)_i&i\in\underline{\mathbb{K}}_7\\
u_i&+\tau\gamma&&-\textcolor{green}{t_i}&&=\underline{\delta}&i\in\underline{\mathbb{K}}_0\\
u_i&+\tau\gamma&&-\textcolor{red}{t_i}&&=\underline{\delta}&i\in\underline{\mathbb{K}}_1\\
u_i&+\tau\gamma&&+\textcolor{red}{t_i}&&=\underline{\delta}&i\in\underline{\mathbb{K}}_2\\
u_i&+\tau\gamma&&+\textcolor{green}{t_i}&&=\underline{\delta}&i\in\underline{\mathbb{K}}_3\\
u_i&+\tau\gamma&&+\textcolor{green}{t_i}&&=\underline{\delta}&i\in\underline{\mathbb{K}}_4\\
&\tau\gamma&+s_i&&+\textcolor{red}{w_i}&=\underline{\delta}&i\in\underline{\mathbb{K}}_5\\
&\tau\gamma&+s_i&&+\textcolor{green}{w_i}&=\underline{\delta}&i\in\underline{\mathbb{K}}_6\\
&\tau\gamma&+s_i&&+\textcolor{green}{w_i}&=\underline{\delta}&i\in\underline{\mathbb{K}}_7\\
&&&&&\\
u\geq0&\gamma\geq0&s\geq0&t\geq0&w\geq0&.\\
\end{array} 
\]
One can see that the constraint matrix is full rank, with $rank = 2n$.\\\\

\begin{algorithm}[H]
        \caption{Escaping $\gamma=0$ }
        \label{alg:gamma0}
        \begin{algorithmic}[1] % [1] adds line numbering
	 \State Inputs $(\underline{\delta}, \underline{x})$.
	  \State Perform Gauss-Jordan elimination with green and red as basis.
            \State Apply one Simplex method change of basis step.
	 \If {there is no row with a nonnegative value on the $\gamma$ column and  a 0 value on the right hand side}
	\State  We've found a point with $\gamma>0$, and it is simplex optimal.
	\State $\bar{x}:=\underline{M}u$, and $\bar{\delta}:=\underline{\delta}-\tau\gamma$.
	\Else
	\State  For each concerned row, do as follow, 
		\[
		\begin{array}{l}
		\text{if}\;i\in\underline{\mathbb{K}}_1\Rightarrow\;\underline{\mathbb{K}}_1:=\underline{\mathbb{K}}_1\setminus\left\{i\right\},\;\underline{\mathbb{K}}_2:=\underline{\mathbb{K}}_2\cup\left\{i\right\}\\
		\text{if}\;i\in\underline{\mathbb{K}}_2\Rightarrow\;\underline{\mathbb{K}}_2:=\underline{\mathbb{K}}_2\setminus\left\{i\right\},\;\underline{\mathbb{K}}_1:=\underline{\mathbb{K}}_1\cup\left\{i\right\}\\
		\text{if}\;i\in\underline{\mathbb{K}}_4\Rightarrow\;\underline{\mathbb{K}}_4:=\underline{\mathbb{K}}_4\setminus\left\{i\right\},\;\underline{\mathbb{K}}_5:=\underline{\mathbb{K}}_5\cup\left\{i\right\}\\
		\text{if}\;i\in\underline{\mathbb{K}}_5\Rightarrow\;\underline{\mathbb{K}}_5:=\underline{\mathbb{K}}_5\setminus\left\{i\right\},\;\underline{\mathbb{K}}_4:=\underline{\mathbb{K}}_4\cup\left\{i\right\}\\
		\end{array} 
		\]
	\State and go to 1.
	\EndIf
	\State Return $(\bar{\delta}, \bar{x})$.
        \end{algorithmic}
    \end{algorithm}

\begin{thm}
\begin{enumerate}
\item At most two Gauss-Jordan eliminations  and two Simplex Method change of basis steps are performed to obtain $\gamma>0$.
\item $x^{\beta}(\delta)$ is a piecewise linear function with respect to $\delta\geq 0$.
\end{enumerate}
\end{thm}
\begin{proof}
\begin{enumerate}
\item ...
\item ...
\end{enumerate}
\end{proof}
\subsubsection{Entire path direct solutions}
This is a first approach that allow entire path solutions computation using Gauss-Jordan eliminations and Simplex Method. LARS like algorithm is also under progress!\\\\
Recall that 
\begin{itemize}
\item $\underline{x}:=0$ is a solution for  $(\mathbb{Q}^{\beta,\delta})$, 
for all $\delta\geq \delta^0:=\ddfrac{1}{\beta}\left\|A^Tb\right\|_{\infty}$.
\item $\underline{x}:=x^{Ridge^{\beta}}$ is a solution for  $(\mathbb{Q}^{\beta,\delta})$, 
\[\text{for all} \;\delta\leq\delta^{Ridge^{\beta}}:=
\underset{x^{Ridge^{\beta}}_i\neq 0}{\min}\left|x^{Ridge^{\beta}}_i\right|.\]
\end{itemize}
\begin{algorithm}[H]
        \caption{Full path computation $(\mathbb{Q}^{\beta,\delta})$, $\delta\geq 0$}
        \label{alg:full0}
        \begin{algorithmic}[1] % [1] adds line numbering
            \State  $\underline{\delta}:=\delta^0$ and $\underline{x}:=0$, $\tau:=+1$,
          \State $L := \left\{(\underline{\delta}, \underline{x})\right\}$
	 \While{ $\underline{\delta} > \underset{i=1,...,n}{\min}\left|\underline{x}_i\right|$}
		\State $(\underline{\delta}, \underline{x})$ := output of Algorithm~\ref{alg:gamma0},
                     \State $L := L\cup\left\{(\underline{\delta}, \underline{x})\right\}$.
	\EndWhile
            \State \Return $\hat{x}$
        \end{algorithmic}
    \end{algorithm}

\begin{algorithm}[H]
        \caption{Full path computation $(\mathbb{Q}^{\beta,\delta})$, $\delta\geq 0$}
        \label{alg:fullRidge}
        \begin{algorithmic}[1] % [1] adds line numbering
            \State  $\underline{\delta}:=\delta^{Ridge^{\beta}}$ and $\underline{x}:=x^{Ridge^{\beta}}$, $\tau:=-1$,
          \State $L := \left\{(\underline{\delta}, \underline{x})\right\}$
	 \While{ $\underline{x}\neq 0$}
		\State $(\underline{\delta}, \underline{x})$ := output of Algorithm~\ref{alg:gamma0},
                     \State $L := L\cup\left\{(\underline{\delta}, \underline{x})\right\}$.
	\EndWhile
            \State \Return $\hat{x}$
        \end{algorithmic}
    \end{algorithm}

We also implemented a parallel version using both Algorithm~\ref{alg:full0} and Algorithm~\ref{alg:fullRidge}, with a critical common section.
\subsection{Cyclic coordinate descent for solving $(\mathbb{Q}^{\beta,\delta})$}
\begin{algorithm}[H]
        \caption{A cyclic coordinate descent for $(\mathbb{Q}^{\beta,\delta})$}
        \label{alg:cda}
        \begin{algorithmic}[1] % [1] adds line numbering
	  \State Precompute $z_j := \sum_{i}^{m}A_{i,j}^2 =\left\|A_{.,j}\right\|^2_2$
            \State Initialize $\hat{x} := 0$, (or something else)
	 \While{ not converged}
		\For{$j=1, ..., n$}
			\State Compute $\alpha_j := \sum_{i}^{m}A_{i,j}(b_i - \sum_{k\neq j} \hat{x}_k A_{i,k}),$
			\State  and set 
\[
\hat{x}_j := 
\left\{ \begin{array}{ll}
0& \left|\alpha_j\right|\leq \delta\beta,\\
\ddfrac{\alpha_j - \delta\beta\text{sign}(\alpha_j)}{ z_j} & \delta\beta <\left|\alpha_j\right|\leq \delta(\beta +z_j),\\
\ddfrac{\alpha_j}{\beta + z_j} &\delta(\beta +z_j) \leq \left|\alpha_j\right|.
\end{array}
\right.
\]
           	 \EndFor
	\EndWhile
            \State \Return $\hat{x}$
        \end{algorithmic}
    \end{algorithm}
\subsection{Interior point method}
We will consider the convex quadratic form of $(\mathbb{Q}^{\beta,\delta})$, see Le Thi et al.

\[
	\left\{\begin{array}{ll}
\min & \left\|Ax-b\right\|^2_2 + \beta\left\| s + \delta v\right\|_2^2 - 2\beta\delta^2 e^Tv\\
		s.t.&\;-s_i\leq x_i\leq s_i,\;0\leq v_i, i=1,...n.
		\end{array} \right.
		\]
Logarithmic barrier for the bounds
\[
\Phi(x, s, v) = -\sum_{i}^{n}log(s_i+x_i)-\sum_{i}^{n}log(s_i-x_i)-\sum_{i}^{n}log(v_i).
\]
\[
\text{dom}\;\Phi = \left\{\left(\begin{array}{l}
x\\
s\\
v\\
\end{array}\right) : -s_i< x_i< s_i,\; 0<v_i,\; i=1,...n.
		     \right\}.
\]
The central path consists of the unique minimizer $(x(t),s(t), v(t))$ of the convex function 
\[
\phi_t(x, s, v) =t\left[ \left\|Ax-b\right\|^2_2 + \beta\left\| s + \delta v\right\|_2^2 - 2\beta\delta^2 e^Tv\right] + \Phi(x, s, v).
\]

Newton method is used to minimize $\phi_t$, i.e, the search direction is computed as the exact solution to Newton system
\[
H\left(\begin{array}{l}
\Delta x\\
\Delta s\\
\Delta v\\
\end{array}\right) = -g.
\]
The Hessian 
\[
H =\nabla^2\phi_t(x, s, v),
\]
and the gradient 
\[
g =\nabla\phi_t(x, s, v),
\]
at the current iterate $(x, s, v)$.\\
...
\section{Sparse Principal Component}
Consider a covariance matrix $Q$.\\\\ 
Writing $Q = A^TA$ and $b := A\alpha$, we can see that 
\[\left\|Ax-b\right\|^2_2=(x-\alpha)^TQ(x-\alpha),
\]
\[
A^T(Ax-b)=Q(x-\alpha)
\]
and 
\[
A^Tb = Q\alpha.
\]
A first approach is to replace ELASTICNET engine by  $(\mathbb{Q}^{\beta,\delta})$ in SPCA algorithm from Zou et al..
\section*{References}
\begin{thebibliography}{99} % The '9' handles up to 9 references

\bibitem{Lethi2016}
Le Thi, H. A. L., Pham Dinh, T., \& Thiao, M.: \textit{Efficient approaches for l2-l0 regularization
and applications to feature selection in SVM.} Applied Intelligence (2016), 45:549–565.

\bibitem{Zou2006}
Zou, H., Hastie, T., \& Tibshirani, R.: \textit{Sparse Principal Component Analysis.} Journal of Computational and Graphical Statistics (2006), 15(2), 265–286.

\bibitem{Friedman2007}
Friedman, J., Hastie, T., Hofling, H., \& Tibshirani, R.: \textit{Pathwise coordinate optimization.} In Annals
of Applied Statistics  (2007).

\bibitem{Kim2007}
Kim, S.J., Koh, K., Lustig, M., Boyd, S. \& Gorinevsky. D.: \textit{An Interior-Point Method for Large-Scale l1-Regularized Least Squares.} IEEE Journal on Selected Topics in Signal Processing (2007), 1(4):606-617.

\end{thebibliography}
\end{document} % This is the end of the document
